%Plantilla basada en "Template for Masters / Doctoral Thesis" (plantilla disponible en writeLaTex) que subió LaTeXTemplates.com

\documentclass[11pt]{book}
\usepackage[paperwidth=17cm, paperheight=22.5cm, bottom=2.5cm, right=2.5cm]{geometry}
\usepackage{amssymb,amsmath,amsthm} %paquete para símbolo matemáticos
\usepackage[spanish]{babel}
\usepackage[utf8]{inputenc} %Paquete para escribir acentos y otros símbolos directamente
\usepackage{enumerate}
\usepackage{graphicx}
%\usepackage{subfig} %para poner subfiguras
\graphicspath{{Img/}} %En qué carpeta están las imágenes
\usepackage[nottoc]{tocbibind}
\usepackage[pdftex,
            pdfauthor={Ilan Jinich Fainsod},
            pdftitle={El problema de admisión a universidades con cotas inferiores},
            pdfsubject={ÁREA DE LA TESIS},
            pdfkeywords={PALABRAS CLAVE},
            pdfproducer={Latex con hyperref},
            pdfcreator={pdflatex}]{hyperref}



\begin{document}

%----------------------------------------------------------------------------------------
%	COMANDOS PERSONALIZADOS
%----------------------------------------------------------------------------------------

%SI TU TESIS TIENE TEOREMAS Y DEMOSTRACIONES, PUEDES DESCOMENTAR Y USAR LOS SIGUIENTES COMANDOS

\renewcommand{\proofname}{Demostración}
%\providecommand{\norm}[1]{\lVert#1\rVert} %Provee el comando para producir una norma.
%\providecommand{\innp}[1]{\langle#1\rangle} 
%\newcommand{\seno}{\mathrm{sen}}
%\newcommand{\diff}{\mathrm{d}}

\newtheorem{teo}{Teorema}[section] 
\newtheorem{cor}[teo]{Corolario}
\newtheorem{lem}[teo]{Lema}
\renewcommand{\qedsymbol}{$\blacksquare$}

\theoremstyle{definition}
\newtheorem{dfn}[teo]{Definición}

\theoremstyle{remark}
\newtheorem{obs}[teo]{Observación}

\newtheorem{eje}{Ejemplo}[section]

\allowdisplaybreaks


%----------------------------------------------------------------------------------------
%	PORTADA
%----------------------------------------------------------------------------------------

\title{TÍTULO DE LA TESIS} %Con este nombre se guardará el proyecto en writeLaTex

\begin{titlepage}
\begin{center}

\textsc{\Large Instituto Tecnológico Autónomo de México}\\[4em]

%Figura
\begin{figure}[h]
\begin{center}
\includegraphics{logo-ITAM_ch.jpg}
\end{center}
\end{figure}

\vspace{4em}

\textsc{\huge \textbf{El problema de admisión a universidades con cotas inferiores y comunes}}\\[4em]

\textsc{\large Tesis}\\[1em]

\textsc{que para obtener el título de}\\[1em]

\textsc{LICENCIADO EN MATEMÁTICAS APLICADAS}\\[1em]

\textsc{presenta}\\[1em]

\textsc{\Large Ilan Jinich Fainsod}\\[1em]

\textsc{\large Asesor: DR RAMON ESPINOSA ARMENTA }

\end{center}

\vspace*{\fill}
\textsc{México, D.F. \hspace*{\fill} 2018}

\end{titlepage}


%----------------------------------------------------------------------------------------
%	DECLARACIÓN
%----------------------------------------------------------------------------------------

\thispagestyle{empty}
\vspace*{\fill}
\begingroup
``Con fundamento en los artículos 21 y 27 de la Ley Federal del Derecho de Autor y como titular de los derechos moral y patrimonial de la obra titulada ``\textbf{TÍTULO DE LA TESIS}'', otorgo de manera gratuita y permanente al Instituto Tecnológico Autónomo de México y a la Biblioteca Raúl Bailléres Jr., la autorización para que fijen la obra en cualquier medio, incluido el electrónico, y la divulguen entre sus usuarios, profesores, estudiantes o terceras personas, sin que pueda percibir por tal divulgación una contraprestación''.

\centering

\hspace{3em}

\textsc{AUTOR}

\vspace{5em}

\rule[1em]{20em}{0.5pt} % Línea para la fecha

\textsc{Fecha}
 
\vspace{8em}

\rule[1em]{20em}{0.5pt} % Línea para la firma

\textsc{Firma}

\endgroup
\vspace*{\fill}


%----------------------------------------------------------------------------------------
%	DEDICATORIA
%----------------------------------------------------------------------------------------

\pagestyle{empty}
\frontmatter

\chapter*{}
\begin{flushright}
\textit{``To life, to life, L'Chaim!''}
%\textit{DEDICATORIA}
\end{flushright}


%----------------------------------------------------------------------------------------
%	AGRADECIMIENTOS
%----------------------------------------------------------------------------------------

\chapter*{Agradecimientos}
%\markboth{AGRADECIMIENTOS23}{AGRADECIMIENTOS} % encabezado 

¡Muchas gracias a todos!


%----------------------------------------------------------------------------------------
%	PREFACIO
%----------------------------------------------------------------------------------------

\chapter*{Prefacio}

\pagestyle{plain}
%\markboth{PREFACIO23}{PREFACIO} % encabezado 

PUEDEN QUITAR ESTA PARTE


%----------------------------------------------------------------------------------------
%	TABLA DE CONTENIDOS
%---------------------------------------------------------------------------------------

\tableofcontents
\newpage


%----------------------------------------------------------------------------------------
%	TESIS
%----------------------------------------------------------------------------------------
\mainmatter %empieza la numeración de las páginas
\pagestyle{headings}

%  Incluye los capítulos en el folder de capítulos

\chapter{Introducción}
\begin{flushright}
\textit{``Matchmaker, matchmaker,
Make me a match, Find me a find,
Catch me a catch''}
\end{flushright}
\section{El problema de admisión a universidades}
Descripción del problema



\begin{dfn}{(Gale Shapley 1962)}
\label{Estable}
Decimos que una asignación de solicitantes a universidades es inestable si existen dos solicitantes $\alpha$ y $\beta$ asignados a  universidades $A$ y $B$ respectivamente, con la propiedad de que $\alpha$ prefiere estar en $B$ que en $A$ y $B$ prefiere tener a $\alpha$ que a $B$. Es decir, existen una universidad y un solicitante que se prefieren entre ellos a sus respectivas asignaciones. Alternativamente una asignación de solicitantes a universidades es estable si no es inestable.
\end{dfn}

\begin{dfn}{(Gale Shapley 1962)}
\label{optima}
Una asignación es considera optima si cada solicitante esta mejor o igual respecto a sus preferencias que en cualquier otra asignación estable. 
\end{dfn}



%DEscripcion
%Definiciones
%Chane esto jala mejor como capitulo y no seccion
%COmentario que explique como abajo es caso particular

\section{El problema del matrimonio estable}
\begin{flushright}
\textit{``Even the worst husband, God forbid, is better than no husband, God forbid.''}
\end{flushright}
%Definicion del problema
%Definiciones alternativas
\begin{dfn}
Definimos la matriz de preferencias para un problema con $n$ y con $m$ mujeres como una matriz con $n$ filas y $m$ columnas donde donde la primera entrada de la celda $(i,j)$ represente el orden de preferencia que le asigna el hombre $i$ a la mujer $j$ y análogamente la segunda entrada representa el orden de preferencia que le da la mujer $j$ al hombre $i$.
\end{dfn}
\begin{eje}
\label{ejemplo matrimonio 1}
Supongamos que contamos con tres hombres $\alpha$, $\beta$ y $\gamma$ y con tres mujeres $A$, $B$ y $C$ y la matriz de preferencias
$$\begin{pmatrix}
& A & B & C \\
\alpha & 1,3 & 2,2 & 3,1 \\
\beta & 3,1 & 1,3 & 2,2 \\
\gamma & 2,2 & 3,1 & 1,3 
\end{pmatrix}$$
aquí por ejemplo el orden preferencias de $\alpha$ es $(A,B,C)$ y el orden de preferencias de $C$ es $(\gamma, \beta, \alpha)$.
\end{eje}

\begin{dfn}
Análogo a la definición \ref{Estable},
decimos que un conjunto de matrimonios es inestable si existe un hombre y una mujer que se prefieren entre si a sus respectivas parejas. Alternativamente un conjunto de matrimonios es estable si no es inestable. 
\end{dfn}

%Ejemplos
\begin{eje}
Retomando el ejemplo \ref{ejemplo matrimonio 1}. \\
El conjunto de parejas $(\alpha, A), \; (\beta, B)\; y\; (\gamma, C)$ es estable porque a pesar de las preferencias de las mujeres cada hombre esta con su primera opción, es decir, los hombres prefieren a su pareja a cualquier otra mujer.\\ 

\end{eje}
\begin{eje}
Retomando el ejemplo \ref{ejemplo matrimonio 1}. \\
El conjunto de parejas $(\alpha, A)$, $(\gamma, B)$ y $(\beta, C)$ es inestable porque $\gamma$ prefiere a $A$ que a su pareja actual ($B$) y A prefiere a $\gamma$ que a su pareja actual $(\alpha)$.
\end{eje}
\subsection{Algoritmo de Gale Shapley}
\begin{eje}
Ejemplo de 4x4
\end{eje}

\begin{teo}[Teorema de Gale Shaley]
\label{teorema de Gale Shapley}
El algoritmo de Gale Shapley termina en una emparejamiento estable.
\end{teo}
\begin{proof}

\end{proof}

\begin{lem}
\label{lema 1}
Bajo el emparejamiento obtenido por Gale Shapley, solo un hombre puede terminar con la última mujer de su lista como pareja. 
\end{lem}

\begin{proof}
Supongamos que en el emparejamiento de Gale Shapley $m$ ($m\geq2$) hombres terminan con la última mujer de su lista como pareja, eso significa que cada uno de esos $m$ hombres invito a salir a todas las mujeres. Entonces cada mujer fue invita a salir por lo menos $m$ veces, lo cual es una contradicción porque el algoritmo acaba cuando invitan a salir a la última mujer y a esta solo la invitan a salir una vez. Por lo tanto, solo un hombre puede terminar con la última mujer de su lista como pareja. 
\end{proof}

\begin{cor}
Si en un emparejamiento estable por lo menos dos hombres están emparejados con la última mujer de sus respectivas listas, entonces existe dos o más emparejamientos estables en el problema.
\end{cor}

\begin{proof}
Llamemos $m$ al emparejamiento estable donde por lo menos dos hombres están emparejados con la última mujer de sus respectivas listas y llamemos $m'$ al emparejamiento de Gale Shapley.
Por el teorema \ref{teorema de Gale Shapley} sabemos que el algoritmo de Gale Shapley siempre acaba en un emparejamiento estable y además en ese emparejamiento estable solo un hombre puede terminar con la última mujer de su lista como pareja por el lema \ref{lema 1}.
Por lo tanto $m$ y $m'$ son diferentes y como ambos son emparejamientos estables entonces el número de emparejamientos estables en el problema es mayor o igual a dos. 
\end{proof}

\begin{cor}
El número máximo de propuestas en el algoritmo es $n^2-n+1$
\end{cor}

\begin{proof}
Por el lema \ref{lema 1} sabemos que a lo más un hombre acaba emparejado con la última mujer de su lista, por lo tanto el peor emparejamiento posible para el algoritmo es uno donde $n-1$ hombres terminan con la penúltima mujer de sus respectivas listas y un hombre termina con la última. Para llegar a esta situación de acuerdo al algoritmo, los $n-1$ hombres deben de realizar $n-1$ propuestas cada uno y el otro hombre debe de realizar $n$ propuestas. Esto es $(n-1)(n-1)+n$ propuestas que es igual a $n^2-n+1$ propuestas.
\end{proof}

\begin{obs}
La complejidad del algoritmo de Gale Shapley es del orden de $n^2$.
\end{obs}

\begin{eje}
Aquí va el ejemplo de 2x2 del lema
Considere la matriz de preferencias 
$$\begin{pmatrix}
&A&B\\
\alpha &1,1& 2,1 \\
\beta & 1,2 & 2,2
\end{pmatrix}$$
\end{eje}

\begin{eje}
Aquí va el ejemplo de 3x3 del lema
\\ Considere la matriz de preferencias $$\begin{pmatrix}
& A & B & C \\
\alpha & 1,2 & 3,2 & 2,1 \\
\beta & 2,3 & 3,1 & 1,2 \\
\gamma & 2,1 & 3,3 & 1,3 
\end{pmatrix}$$
\end{eje}



\thispagestyle{empty}
\chapter{Problemas NP-Completos}

\section{Clases de complejidad}
Algo que nos interesa mucho cuando atacamos un problema de computación es que tan dificíl es de resolver y si existe una mejor manera de hacerlo. La complejidad computacional nos permite estudiar esto. Para encarrilarnos hay que empezar con algunas definiciones.
\begin{dfn}
Se dice que un problema es un problema de decisón si admite dos posibles respuestas ``Si'' o ``No''. Resolverlo es encontrar todas las soluciones en las que la respuesta es ``Si''.
\end{dfn}
\begin{dfn}
Un problema esta en la clase NP si existe un algoritmo polinomial para verificar que la respuesta de una solución dada es ``Si''.
\end{dfn}
\begin{eje}[Problema del clima]
Supongamos que Mariana le dice a Adrián que esta lloviendo, si Adrián decide verificar si esto es cierto solo tiene que realizar tres pasos:
\begin{enumerate}
\item Caminar a la ventana.
\item Abrir la ventana.
\item Ver el cielo.
\end{enumerate}
Por lo tanto a Adrián le toma un tiempo constante verificar si esta lloviendo lo que implica que este problem esta en NP.
\end{eje}
\begin{dfn}
Se dice que un problema de decisión esta en P si existe un algoritmo polinomial que lo resuelve.
\end{dfn}
\begin{obs}[$P  \subseteq NP$]
Si un problema de decisión pertenece a P entonces pertenece a NP.
\end{obs}
%\begin{proof}
%Supongamos que $A$ es un problema de decision en P y $b$ es una solución dada de $A$, podemos verificar si la respuesta de la solución es ``Si'' resolviendo el problema en tiempo polinomial.
%\end{proof}
\begin{obs}
Si A es un problema de decisión en P y existe un algoritmo polinomial que reduce resolver el problema de decisión B a resolver el problema A entonces B esta en P.
\end{obs}

Vale la pena mencionar algunos ejemplos de problemas que estan en P:
\begin{enumerate}
\item El problema de admisión a universidades.
\item Verificar si un numero es primo.
\item Encontrar la ruta más corta entre dos vertices en una gráfica.
\item El problema de la ruta critica. 
\end{enumerate}

Una conjetura famosa que vale la pena mencionar es si $P=NP$, es decir si todo problema en NP tambien se enecuentra en P. La persona que lo demuestra además de ganar fama y traajo es acredor a un permio de un millon de dolares. 

\section{El $k-SAT$}
\section{Problemas NP-Completos}
\begin{flushright}
\textit{``Chaos is a ladder'}
\end{flushright}
\begin{dfn}
Qué es un problema NP-completo. 
\end{dfn}
\begin{teo}[Teorema de Cook-Levin]
El 3-SAT es NP-Completo.
\end{teo}
Donde esta la demostración. 
\thispagestyle{empty}
\chapter{Conclusiones}

Concluyo que
\thispagestyle{empty}


%----------------------------------------------------------------------------------------
%	APÉNDICES
%----------------------------------------------------------------------------------------

\addtocontents{toc}{\vspace{2em}} % Agrega espacios en la toc

\appendix % Los siguientes capítulos son apéndices

%  Incluye los apéndices en el folder de apéndices

\chapter{Álgebra Lineal}

Cosas.
\thispagestyle{empty}
%\include{Apendices/AppendixB}
%\include{Apendices/AppendixC}

\addtocontents{toc}{\vspace{2em}} % Agrega espacio en la toc


%----------------------------------------------------------------------------------------
%	BIBLIOGRAFÍA
%----------------------------------------------------------------------------------------
\backmatter
\nocite{*}
\bibliographystyle{plain}
\bibliography{bibliografía.bib} %Aquí ponen el nombre del archivo .bib






\end{document}